\section{Discussion, Limitations, and Ethics}
\label{sec:discussion}

\looseness=-1
The strongest version of this paper is not an argument against safety
alignment. It is an argument for evaluating where safety controls belong in an
agentic system. Across the cross-family evaluation, visible refusal is only one
small part of the behavior. In Gemma, evidence grounding separates the
conditions after both models read relevant files; in Qwen, the comparison is
mostly flat; in Llama, tool-interface reliability dominates the abliterated
condition. No family passes the hard-validation checks.

\looseness=-1
The main methodological implication is that autonomous security-agent
evaluations should not collapse success, refusal, unsafe action, and evidence
grounding into one score. These outcomes answer different questions. A refusal
label asks whether the model declined an authorized step; a success predicate
asks whether the final artifact satisfies the task; an unsafe-action label asks
whether the sandbox policy was violated; and grounding asks whether the trace
and artifact contain the local files, symbols, issue terms, and evidence needed
to justify the claim. The Gemma result appears in this separation: the aligned
model is not simply refusing, yet it is less often producing the concrete
defensive evidence required by the task.

\looseness=-1
The blind audit is therefore a check on the interpretation, not an alternate
success judge. The deterministic grounding score is intentionally mechanical,
so the audit asks whether independent reviewers see the same coarse structure
in redacted traces when condition labels, original failure labels, checker
reasons, and grounding scores are hidden. Inter-reviewer agreement is high for
success and refusal judgments but low for hard-validation subtypes, and each
reviewer agrees less with the lexical success labels (Section~\ref{sec:results}).

\looseness=-1
The deployment lesson is also narrower than ``remove alignment.'' The safer
interpretation is that policy should be evaluated at the system boundary where
authorization, tool permissions, trace review, redaction, and final evidence
checks can all be observed. For an open-weight model, whoever runs the model
writes the prompt, so prompt text cannot establish authorization; the
surrounding system must.

\subsection{Limitations}
\label{sec:limitations}

\looseness=-1
\textbf{Derivatives are probes, not counterfactuals.} Each model pair uses one
less-restricted derivative with its own refusal-removal procedure (Section~\ref{sec:background}). We did not vary the abliteration
method or strength, so effects are attributable only to these derivatives;
comparing methods and strengths is left to future work.
\textbf{Capability confound.} The Gemma derivatives also do better on
non-security controls, the aligned 26B model mostly fails through repeated tool
calls that the repair controller does not cover (129/150 runs), and the Llama
derivative is much worse at following the
tool protocol, so capability and interface reliability vary with alignment
condition. A causal claim about alignment, which we do not make, would require
pairs matched on non-security performance. \textbf{Controller.} The repair controller changes effect
magnitudes (Section~\ref{sec:results}); its deterministic rules are logged per
trace (Appendix~\ref{app:controller-repairs}).
\textbf{Scale and measurement.} Absolute success is low, no condition passes the
lexical proof-of-trigger or patch-verification predicates, and the 30 tasks are
stratified across five classes rather than sampled from security work at large.
Success and grounding are both lexical substring checks, and the trace audit
uses LLM reviewers rather than human annotators, with low agreement on
hard-validation subtypes (\(\kappa=0.21\)).
\textbf{Scope.} The authorization-prompt ablation covers Gemma 4 31B only, and
all security and control runs use one backend (\texttt{llama.cpp}), 4-bit GGUF quantizations, one
sampling configuration, and five unseeded replicates per task, without network
isolation.

\subsection{Ethics and Dual Use}
\label{sec:ethics}

\looseness=-1
This is dual-use research. All experiments use original local tasks with known
defects, scoped tools, and no third-party targets. Released traces are redacted
and audited, and no jailbreak prompt was applied in any reported run. The less-restricted derivatives are not safe
defaults: two 26B derivative runs installed packages from the network
(Appendix~\ref{app:out-of-scope}), and we do not endorse deploying them. Production systems still
need alignment, sandboxing, policy checks, and evidence verification.
